\chapter{Registros, tablas extendidas y entrega digital}
\section{Registros UART}
Para radiografías se conservará una fila por imagen con identificador, clase real, clase predicha, puntuaciones, tiempo de inferencia y estado de comunicación. Para HIL se conservará una fila por muestra con tiempo, estados, error angular, acción, ciclos, inferencia, RTT e incumplimientos de plazo.

\section{Resultados adicionales}
Este anexo contendrá tablas completas por semilla, condición inicial, perturbación y variante. Las figuras de la memoria se generarán desde estos registros. Se incluirán las trayectorias fallidas además de ejemplos satisfactorios, de modo que las conclusiones no dependan sólo de una selección visual favorable.
\pendiente{Agregar tablas extendidas y manifestar qué archivos originan cada figura del capítulo de resultados.}

\section{Estructura digital solicitada por la norma}
La normativa institucional solicita carpetas \texttt{ADOBE}, \texttt{IMAGENES}, \texttt{SOFTWARE}, \texttt{ANEXOS} y \texttt{AUDIOVISUAL}. La primera contiene el informe PDF; las siguientes agrupan material gráfico, programas, antecedentes y videos. El proyecto Overleaf corresponde a la fuente editable y no reemplaza por sí solo esa organización de entrega.

La norma también describe ejemplares y carátulas para CD. Si la unidad académica aplica actualmente un procedimiento digital diferente, deberá confirmarse al momento de entregar. La plantilla sigue el documento proporcionado y no presupone una actualización administrativa.
